% ---------------------------------------------------------------------------
% Fast preview of intro_icra.tex ALONE, in the real ICRA class.
% Not the submission file --- that is root.tex. This exists so the Introduction
% can be iterated on without rebuilding the whole paper.
%
%   bash build.sh            one-shot build -> preview.pdf
%   bash build.sh --watch    rebuild whenever intro_icra.tex changes
%   bash build.sh root       build the real paper (root.tex) instead
% ---------------------------------------------------------------------------
\documentclass[letterpaper, 10pt, conference]{ieeeconf}
\IEEEoverridecommandlockouts
\overrideIEEEmargins

\usepackage{graphics}
\usepackage{mathptmx}
\usepackage{times}
\usepackage{amsmath}
\usepackage{amssymb}

\title{\LARGE \bf Sim-to-Real Gentle Manipulation of Food}
\author{PREVIEW BUILD --- Introduction + Related Work}

\begin{document}
\maketitle
\thispagestyle{empty}
\pagestyle{empty}

% Placeholder for Fig. 1 (the coverage-failure map, PAPER_TODO A4). Keeps the
% \ref resolving and shows how much column space the real figure will consume.
\begin{figure}[t]
  \centering
  \framebox[0.92\columnwidth]{\rule{0pt}{3.4cm}\parbox{0.8\columnwidth}{\centering
    \textbf{Fig.~1 placeholder}\\[2pt]
    \footnotesize Coverage-failure map: hardware failures of the 50-real-demo
    policy vs.\ (workspace location $\times$ object size), misgrasp and
    over/under-squeeze as distinct markers. \\ \texttt{PAPER\_TODO A4}}}
  \caption{Placeholder --- replace with the real coverage-failure map.}
  \label{fig:coverage}
\end{figure}

% ---------------------------------------------------------------------------
% INTRODUCTION --- v4 (2026-08-26). Register calibrated against DextrAH-G's
% actual introduction (arXiv:2407.02274v2), fetched verbatim.
%
% CHANGE FROM v3 (important, do not undo):
%   Stiffness is NO LONGER a pillar of the argument. The policy observes only a
%   point cloud, so it cannot be stiffness-ADAPTIVE by construction; leading on
%   stiffness would advertise a problem we do not solve. Our own failure data
%   implicates POSE and SIZE, not stiffness. Stiffness now appears once, in P4,
%   as honest ROBUSTNESS (material randomisation -> a grip that works across the
%   range, not one tuned to a value). Prepared answer if a reviewer presses:
%   "the policy is not stiffness-adaptive and cannot be; randomising material
%   during demonstration generation yields robustness across the range. Whether
%   stiffness-adaptive behaviour would help is open and we do not address it."
%
%   2026-08-26: contribution 2 further softened to DESCRIBE the training
%   ("across randomised object pose, size, shape, and material") rather than
%   CLAIM an outcome ("robust across the material range"). We have no stiffness
%   sweep demonstrating robustness; do not reinstate the outcome wording unless
%   that experiment exists.
%
%   Also restores the Fig. 1 reference (the 50-demo coverage-failure map), which
%   v3 dropped --- it is our own strongest motivating evidence.
%
% Structure: P1 problem | P2 prior work + what it leaves | P3 coverage evidence
% (cited) + our own failure map | P4 proposal (three beats only: FEM stress-aware
% scripted demonstrator -> many categories with within-category randomisation ->
% one policy, deployed real) | P5 contributions, standalone.
%
% P4 is deliberately NOT a full pipeline walkthrough (user, 2026-08-26). Keep it
% to the three beats; detail belongs in Method. The specialist/self-distillation
% stage is intentionally absent from the spine --- it is ablated, and leaving it
% out means the paragraph survives either outcome.
% ---------------------------------------------------------------------------
\newcommand{\method}{\textsc{Method}}   % TODO: name the framework

% Pending numbers --- defined once so a landing result is a one-line edit.
% Keep them bold while unresolved so they are impossible to miss in the PDF.
\newcommand{\ndemos}{\textbf{[XX\,k]}}      % synthesised demonstrations
\newcommand{\ncats}{\textbf{[N]}}           % fragile food categories
\newcommand{\nheld}{\textbf{[M]}}           % held-out categories (zero-shot), if reported
\newcommand{\nreal}{\textbf{[55]}}          % real teleop demonstrations in the co-train slice
\newcommand{\headline}{\textbf{[XX\%]}}     % headline real-robot success rate

\section{INTRODUCTION}

Handling delicate food is a skill people exercise without thought, yet it remains difficult for
robots. Automating it matters to food processing, agriculture, and assistive settings, where
produce tears, cracks, and bruises under modest compression and the damage is
irreversible~\cite{zhu2025deformable,blanco2024t,wang2025damageless}. A successful grasp must be
right in two respects at once: where and in what orientation the fingers close, and how far they
close. Closing across a thin edge, or straddling a stem, fails at a width that would have been
correct elsewhere on the same item. Both quantities depend on the item's pose, size, and shape,
and these vary continuously within a single food category and change entirely across categories,
so a policy that handles one item well need not handle its neighbour.

Existing approaches (see Section~\ref{sec:related}) address this largely through hardware and
sensing. Compliant and pneumatic end-effectors bound contact pressure
mechanically~\cite{wang2024softgripper}, and tactile or force-controlled grippers close the loop
on contact at run time, handling delicate items well~\cite{kang2026learning,vtam}. While
effective, these methods change what the robot can \emph{observe} rather than where its
\emph{supervision} comes from: the policies are learned by imitation from human demonstrations
recorded on real hardware, one item and one pose at a time. Stress-aware grasp planning targets
the damaging quantity directly~\cite{pan2020stressmin}, and finite-element simulation can
evaluate grasp outcomes on deformable objects at scale~\cite{defgraspsim}, but neither has been
used to supply the demonstrations a policy trains on. Acquiring broad coverage of the physical
variation therefore remains a human-labour problem.

That coverage is what determines generalisation. Lin \emph{et al.} show that policy performance
follows a power law in the number of environments and objects, and that demonstrations per
object saturate quickly: diversity, not count, is what buys
generalisation~\cite{lin2025datascaling}. Precision sharpens the requirement, since for
high-precision tasks the number of demonstrations needed grows super-exponentially as the
tolerance tightens, against a limit precision that itself depends on the quality of the
expert~\cite{curseofprecision}. We observe the shortfall directly on hardware: a policy trained
on 50 real demonstrations does not degrade uniformly but fails in a structured way, misgrasping
at particular workspace locations and mis-squeezing at particular object sizes
(Fig.~\ref{fig:coverage}). Data-generation systems already multiply a few hundred human
demonstrations into tens of thousands by adapting them to new object
placements~\cite{mandlekar2023mimicgen}. However, such adaptation is geometric and reuses the
grips a human chose; it does not decide afresh how firmly to close on an item it has not seen.

To address these challenges, we propose \method, a framework that generates gentle
demonstrations autonomously in simulation and trains from them a single policy spanning many
fragile food categories. Our key observation is that internal stress, the quantity that
determines whether an item is damaged and which no sensor on our robot reports, is computed
exactly inside a deformable simulator; we therefore treat gentleness as a planning objective
rather than a sensed signal. A finite-element model predicts the stress that closing to a
commanded width induces, and a scripted demonstrator grasps each item at the pose and width
that minimise it, subject to holding the item against gravity. We run this demonstrator across
\ncats{} fragile food categories while randomising object pose, size, shape, and material within
each, and collect \ndemos{} demonstrations with no human teleoperation. A single
category-conditioned policy is trained on the pooled data, co-trained with \nreal{} real
teleoperated episodes~\cite{simrealcotrain}, and deployed on a $7$-DoF arm that observes the
scene as a point cloud from one external depth camera---chosen for its narrow sim-to-real
appearance gap---with no tactile or force sensing in the loop. On hardware it grasps \ncats{}
categories at \headline{} success while handling in-category variation in pose, size, and shape.

The main contributions of this work include: 1) a stress-aware grasp planner built on a
width-controlled, distributed-pad contact model, which selects both where to grasp and how far
to close on a deformable item, 2) an autonomous demonstration generator that produces gentle,
successful grasps across randomised object pose, size, shape, and material with no human
demonstration, 3) a single category-conditioned policy that handles \ncats{} fragile food
categories from a point cloud alone, without tactile or force sensing, and 4) a real-robot
evaluation of that policy across categories and across in-category variation in pose, size, and
shape, a step towards general gentle manipulation of everyday food.


% ---------------------------------------------------------------------------
% RELATED WORK --- v1 (2026-08-26).
%
% Ordered to match the strands the Introduction forward-references, so that
% "(see Section II)" in P2 actually pays off:
%   A. gentle / damage-free handling of fragile objects   (hardware strand)
%   B. contact sensing for delicate manipulation          (sensing strand)
%   C. grasp quality for deformable objects               (stress-aware strand)
%   D. demonstration data: coverage and automatic generation
%   E. sim-to-real for deformable manipulation
%
% Register: ICRA/DextrAH-G. Each subsection states what the line achieves, then
% one sentence on what it leaves open for us. NO dismissals --- every strand is
% credited with what it actually does. See PAPER_TODO "Related-work positioning".
%
% Sources and verification status: docs/paper/related_work.md. Entries marked
% [STUB - FILL IN] in reference.bib are the user's own keys, not yet filled.
% ---------------------------------------------------------------------------
\section{RELATED WORK}
\label{sec:related}

\subsection{Gentle manipulation of fragile objects}
Damage-free handling of fragile produce is a long-standing goal in agricultural and food
automation, and remains open across the pre- and post-harvest
pipeline~\cite{wang2025damageless}. A large part of the effort is mechanical: compliant,
pneumatic, and enveloping end-effectors bound contact pressure by construction, so that a
grasp is gentle whether or not the robot reasons about
it~\cite{wang2024softgripper,zhu2025deformable}. This is effective and complementary to our
work, but compliance is a fixed property of the hardware; it does not choose where on an item
to close, which is one of the two quantities a gentle grasp must get right.

\subsection{Contact sensing for delicate manipulation}
A second line equips the robot to perceive contact as it happens. Kang \emph{et al.} pair a
low-cost force-controlled parallel jaw with a reactive policy and report reliable handling of
chips, tofu, and cherry tomatoes~\cite{kang2026learning}, and video-tactile models attain high
success on contact-rich tasks by fusing tactile streams with
vision~\cite{vtam}. Deformation-based tactile sensing has a known blind spot in the very
regime we target---it requires measurable indentation to produce a signal---which the tactile
community is itself addressing with deformation-independent
designs~\cite{lighttact}. Our work is orthogonal to this line rather than competing with it:
contact feedback supplies run-time adaptation that offline supervision does not, and a
force-controlled gripper trained on stress-planned demonstrations would be a natural
combination. What contact sensing does not change is where supervision comes from; these
policies are trained by imitation from human demonstrations collected on real hardware, and
Kang \emph{et al.} build a dedicated teleoperation device for exactly that purpose.

\subsection{Grasp quality for deformable objects}
Classical grasp metrics score whether an object can be held against external wrenches, which
says little about whether holding it causes damage. Pan \emph{et al.} address this with a
stress-minimising metric that accounts for object material and indicates fragile
regions~\cite{pan2020stressmin}. Their formulation assumes point contacts and general
multi-contact grasps; a two-pad parallel jaw on a smooth rounded object falls outside that
regime, and in our implementation the resistible-wrench hull becomes degenerate, so the metric
does not discriminate among candidate grasps on our objects (Section~\ref{sec:method:synth}).
Finite-element simulation has also been used to characterise grasp outcomes on deformables at
scale: DefGraspSim evaluates parallel-jaw grasps across a large object set and reports stress,
deformation, and stability metrics with good sim-to-real
correspondence~\cite{defgraspsim}. That work \emph{evaluates} grasps; we use the same physics
to \emph{search} for them, and then to generate the demonstrations a policy trains on.

\subsection{Demonstration data: coverage and generation}
How much data an imitation policy needs, and of what kind, has been studied directly. Lin
\emph{et al.} find that generalisation follows a power law in the number of environments and
objects while demonstrations per object saturate quickly, so diversity dominates
count~\cite{lin2025datascaling}; for high-precision tasks the requirement grows
super-exponentially as the tolerance tightens~\cite{curseofprecision}. Since diversity is
expensive to collect by hand, a growing body of work generates it. MimicGen adapts a few
hundred human demonstrations into tens of thousands by transforming them to new object
placements~\cite{mandlekar2023mimicgen}, and related systems extend the idea to dexterous and
dynamic settings. These systems reuse the grasps a human chose and adapt them geometrically.
Our demonstrator instead plans each grasp from a physics objective, so it can select a grip
for an object configuration no human has demonstrated.

\subsection{Sim-to-real for deformable manipulation}
Transferring deformable-object skills from simulation is difficult because contact and
material dynamics are hard to model~\cite{yin2021modeling,arriola2020modeling}, and early work
established the basic feasibility of sim-to-real reinforcement learning in this
setting~\cite{blanco2024t}. Rather than pursue zero-shot transfer, we adopt sim-and-real
co-training, in which a small slice of real demonstrations accompanies a large simulated
set~\cite{simrealcotrain}; we use a point cloud from a single depth camera for the same
reason, its appearance gap between simulation and reality being narrower than that of RGB.


% ---------------------------------------------------------------------------
% METHOD (grasp synthesis + demonstration generation) --- condensed to ~1 column-page.
% Checked against the FROZEN v4.1 recipe actually used for the 12-object collection
% (26 of 28 control fields byte-identical across objects) and against
% docs/grasp_synthesis_model.md (verified-against-code, 2026-08-27).
%
% CORRECTIONS vs the previous draft (do not revert):
%  1. Pad profile is NOT parabolic. Nodes are pushed to a common FLAT plane with a
%     thin smoothstep edge fillet (width 0.12*min(h1,h2)). width_grasp.py:149-155
%     records that the fixed-delta/domed profile was rejected: it "domes a flat face
%     and never flattens a curved one".
%  2. The contact-area floor is NOT a hard constraint in the frozen recipe
%     (grasp_area_min_mm2: auto => area_min=0.0 during search). It is a
%     POOL-RELATIVE POST-SELECTION: restrict to grasps under 0.8*yield, keep the
%     upper half by BOTH worst-pad area and alignment, then take the best score.
%     finger_grasp.py:939-946.
%  3. The camera-azimuth term was missing entirely. It is ACTIVE
%     (cam_azimuth_max_deg: 60) and is a SHAPED PENALTY, not a hard bound
%     (finger_grasp.py:384-386, CAM_AZ_SLOPE=5000/deg of excess; it also bounds
%     the CMA seed fan, :838-847). yaw_max_deg is the hard structural bound.
%     TWO DISTINCT MECHANISMS, do not conflate:
%       - w_occ  (ray-based occlusion FRACTION) = 0.0 in every run, inert in the
%         score, computed for audit only (:436,:542). NOT in the paper.
%       - cam_azimuth_max_deg (azimuth hinge penalty) = ACTIVE. This is the
%         w_az term in Eq. (score). It is a geometric PROXY for occlusion, not
%         a measurement of it --- the wording must not imply otherwise.
%  4. Peak term carries an E factor: w_peak * E * hi_1.
%  5. Added the auto-scaled parameters (the "no per-category constants" claim) and
%     the second width-rescan round.
%  6. Removed the dangling \ref{sec:method:supervision} --- the lead/dwell material
%     was cut from the paper.
%
% MUST NOT CLAIM (model doc warnings): frictional or tangential FEM contact (it is
% normal-only, one scalar constraint per node); force closure or dynamic stability
% (holdability is a quasi-static scalar friction inequality); that occlusion enters
% the objective (w_occ = 0 in every run).
% ---------------------------------------------------------------------------
\section{METHOD}

\subsection{Problem setting}
\label{sec:method:problem}
A robot arm with a parallel-jaw gripper must lift a single fragile object from a table and hold
it without damaging it. Object category, pose, size, shape and material vary between episodes.
Perception is a single fixed external depth camera; there is no tactile or force sensing in the
loop. At each control step the policy observes
$o_t = (\mathbf{p}_t, \mathbf{q}_t, w_t, \mathcal{P}_t)$: end-effector position, orientation as
a unit quaternion, measured gripper width, and an $N_p\!=\!1024$ point cloud in the robot base
frame. It emits an \emph{absolute} target pose and width
$a_t = (\mathbf{p}^{\mathrm{cmd}}_t, \boldsymbol{\theta}^{\mathrm{cmd}}_t, w^{\mathrm{cmd}}_t)
\in \mathbb{R}^{7}$, with orientation as Euler angles. An episode succeeds if the object centre
enters a height band and stays there for $T_{\mathrm{hold}}$ steps. We call a grasp
\emph{gentle} when peak internal von~Mises stress stays below the material's yield threshold;
success is measured on hardware, gentleness in simulation.

\subsection{Stress-aware grasp synthesis}
\label{sec:method:synth}
For each object instance the demonstrator plans a $7$-DoF grasp
$x=(\mathbf{t},\mathbf{r},\omega)$: tool-centre position, orientation, and \emph{commanded}
gripper width. Planning runs on the mesh, outside the physics simulator.

\noindent\textbf{Width-controlled contact model.}
A real parallel jaw is position-controlled: one commands a width and the grip force is an
outcome. We model this directly. The object is tetrahedralised and a linear-elastic stiffness
matrix assembled; since it is free-floating the system is singular in the six rigid-body modes,
so we solve a bordered inertia-relief system, factored once per geometry. Two rigid pads, whose
rectangular footprints are taken from the real gripper's finger meshes, are closed to width
$\omega$. Boundary nodes inside a footprint are pushed onto the pad's flat plane and constrained
\emph{only} along the closing axis, leaving tangential motion free so the material bulges; a
thin smoothstep fillet at the pad edge avoids a flat-punch stress singularity. Contact is
normal-only, and friction enters solely through the holdability test below.

\noindent\textbf{Material scaling.}
Because the displacement is prescribed rather than the force, the deformation field is
independent of Young's modulus, so below yield both stress and grip force are linear in it,
\begin{equation}
  \sigma(E) = E\,\sigma_1, \qquad F(E) = E\,F_1 .
  \label{eq:E-linear}
\end{equation}
One solve at $E\!=\!1$ therefore yields every material by scalar multiplication, which makes
randomising stiffness across the demonstration set essentially free. A grasp is
\emph{holdable} when Coulomb friction at the two pads carries the object's weight with an
acceleration margin, $2\mu F(E) \ge m\,(g+a)$; we use $\mu\!=\!0.7$ and $a\!=\!g$, a quasi-static
$2g$ criterion rather than a wrench-closure analysis.

\noindent\textbf{Objective.}
Let $\sigma_{\mathrm{blk}}$ be the mean of the top decile of element stresses with
contact-adjacent elements masked out, so that it reports \emph{bulk} deformation rather than the
mesh-dependent contact singularity, and let $\sigma_{p98}$ be its unmasked $98$th-percentile
counterpart. Among holdable candidates we maximise
\begin{equation}
  S(x) = -\,\sigma_{\mathrm{blk}}
         - w_{\mathrm{al}}\bigl(1-\mathrm{al}\bigr)
         - w_{\mathrm{pk}}\,E\,\sigma_{p98}
         - w_{\mathrm{pr}}\,\frac{F}{A_{\min}},
  \label{eq:score}
\end{equation}
where $\mathrm{al}$, the mean absolute cosine between the closing axis and the local surface
normals, rejects grazing and edge contacts; the peak and pressure terms penalise \emph{local}
damage that the masked bulk statistic hides, $A_{\min}$ being the area of the worse pad.

Two soft constraints act on $S$ from outside. Candidates failing a feasibility
ladder---table collision, mesh penetration, missing or excessive contact, invalid solve,
unholdable---receive a \emph{shaped} penalty, so the search retains a gradient toward the
feasible band rather than a flat infeasible floor. Separately, a grasp whose closing-axis
azimuth about the camera direction exceeds $60^\circ$ is penalised linearly in the excess, a
geometric proxy for the fingers occluding the object in the very observation the policy must
act on; the same cone bounds the seed fan below.

\noindent\textbf{Search and selection.}
$S$ is optimised by multi-start CMA-ES over $(\mathbf{t},\mathbf{r},\omega)$, seeded from
canonical closing axes and a yaw fan about the camera-perpendicular direction, since a single
start reliably misses the gentle basin. A second pass rescans width over the best distinct
poses, the widest holdable grasp being the gentlest. Because the factorisation is reused across
all candidates sharing a geometry, each evaluation costs one back-substitution, which is what
permits $\sim\!10^3$ evaluations per grasp. Rather than impose a hand-set contact-area floor, we
select pool-relatively: among feasible grasps under $0.8\times$ yield we keep the upper half by
\emph{both} worst-pad area and alignment, then take the best score. This is what separates
enveloping grasps from stem-type ones, which score \emph{better} on bulk stress---squeezing a
stiff, slender part deforms the body less---but differ sharply in contact area. Four remaining
parameters are derived per object rather than tuned: the width cap from the median local
cross-section, the yaw bound from the largest extent, the closure margin from the smallest
extent, and the area criterion above. No per-category constant is set by hand.

\subsection{Demonstration generation}
\label{sec:method:demos}
Planned grasps are executed in a deformable (MPM) simulator, producing
(observation,\,action) trajectories in the same format a real teleoperation session yields.

\noindent\textbf{Randomisation.}
Each episode samples object pose, uniform scale, and a procedural shape deformation, together
with the material parameters that Eq.~\eqref{eq:E-linear} makes free to vary. Each category
additionally carries several meshes representing distinct variants within the type
(Fig.~\ref{fig:meshes}).

\noindent\textbf{Execution.}
Each environment runs an independent phase machine---approach, settle, close, firm, lift,
hold---so environments may progress at different rates while commands are still issued as one
batched step. At the close--firm boundary the measured contact signal is read once, and a grasp
that came out weak receives a small additional closure. Slip during the lift triggers a bounded
regrasp, and the failed attempt is deliberately retained in the demonstration, so recovery
behaviour is represented in the data.

\noindent\textbf{Approach trajectory.}
Rather than interpolating linearly from home to the grasp pose, the approach reproduces the
structure we measured in human teleoperation: lateral alignment completes while the gripper is
still well above the object, and descent continues afterwards. We implement this as per-axis
progress within a single phase, with no waypoint and no stop,
\begin{equation}
  s_{xy}(\alpha) = \mathrm{smoothstep}\!\bigl(\min(\alpha/f,\,1)\bigr), \qquad
  s_{z}(\alpha) = \alpha ,
  \label{eq:approach}
\end{equation}
with $\alpha\in(0,1]$ and $f$ drawn per environment, setting how early the lateral motion
finishes; orientation is slerped over the full $\alpha$. Lateral velocity reaches zero only at
its own arrival, while $z$ is still moving, so the path never dwells. The phase duration is set
from the profile's arc length and a reference speed rather than being fixed, so per-step speed
is approximately constant across episodes; a fixed duration would instead make commanded speed
depend on how far the object happened to spawn.


\begin{figure}[t]\centering
  \framebox[0.9\columnwidth]{\rule{0pt}{2.2cm}\parbox{0.8\columnwidth}{\centering
  \textbf{Fig. placeholder} \\ \footnotesize per-category mesh variants + scale/shape DR}}
  \caption{Placeholder --- mesh variants and shape randomisation.}\label{fig:meshes}
\end{figure}

\bibliographystyle{IEEEtran}
\bibliography{IEEEabrv, reference}

\end{document}
